\chapter{Środowisko uruchomieniowe i weryfikacja działania}\label{ch:testy}

Niniejszy rozdział charakteryzuje środowisko uruchomieniowe wykorzystane do wdrożenia systemu oraz przedstawia metodykę i wyniki weryfikacji działania poszczególnych modułów, obejmującej automatyczny potok CI/CD, izolowane testy jednostkowe w piaskownicy oraz uprzywilejowane testy akceptacyjne.

\section{Środowisko uruchomieniowe i konfiguracja bezpieczeństwa}\label{sec:srodowisko-testowe}

System został wdrożony i przetestowany w dwuwęzłowej architekturze sieciowej łączącej zarządzaną stację roboczą z systemem Microsoft Windows oraz dedykowany serwer kopii zapasowych z systemem Linux. Parametry obu węzłów oraz skonfigurowanych usług zestawiono w tabeli~\ref{tab:srodowisko-testowe}, natomiast całościową topologię połączeń, komponentów i przepływów protokołów komunikacyjnych przedstawiono na rysunku~\ref{fig:environment-topology}.

\Figure[1.0]{environment_topology.png}{Topologia dwuwęzłowego środowiska uruchomieniowego, separacja ról oraz kierunki przepływu protokołów komunikacyjnych}{fig:environment-topology}

\begin{table}[H]
    \centering
    \caption{Konfiguracja węzłów środowiska uruchomieniowego i testowego}\label{tab:srodowisko-testowe}
    \tablefont
    \begin{tabularx}{\linewidth}{L{3.2cm} L{5.4cm} Y}
        \toprule
        \hcell{Parametr} & \hcell{Stacja robocza (Klient)} & \hcell{Serwer kopii zapasowych} \\
        \midrule
        System operacyjny & Microsoft Windows 10 Pro & Ubuntu Linux 26.04 LTS \\
        Adresacja IP (LAN) & \code{10.0.2.20} & \code{10.0.2.30} \\
        Środowisko skryptowe & Windows PowerShell 5.1 (\textit{Desktop Edition}), .NET 4.8 & Powłoka systemowa \cmd{bash}, narzędzia \cmd{flock}, \cmd{sha256sum}, \cmd{base64} \\
        Usługi komunikacyjne & OpenSSH Client + OpenSSH Server (\code{sshd} na porcie TCP 22) & OpenSSH Server (port TCP 22) + SFTPGo (port TCP 2022) \\
        Konta operacyjne & Konto administratora \code{Student} + konto \code{SYSTEM} (\textit{Watcher}) & Konto systemowe \code{fwbackup} + konto usługi \code{sftpgo} \\
        Katalogi chronione & \file{C:\FW}, \file{C:\ProgramData\FirewallAudit}, \file{C:\ProgramData\FirewallBackup} & \file{/etc/fwbackup}, \file{/srv/fwbackups/student/backups} \\
        \bottomrule
    \end{tabularx}
\end{table}

W celu zapewnienia odporności wdrożenia na lokalną eskalację uprawnień oraz nieautoryzowany dostęp sieciowy zastosowano następujące środki utwardzania konfiguracji:
\begin{enumerate}
    \item \textbf{Ochrona katalogów skryptów i dzienników na stacji Windows} - katalogi \file{C:\FW} oraz \file{C:\ProgramData\FirewallAudit} posiadają wyłączone dziedziczenie uprawnień DACL (\cmd{icacls.exe /inheritance:r}) i przyznają prawa zapisu wyłącznie kontom \code{SYSTEM} (\code{S-1-5-18}) oraz \code{Administratorzy} (\code{S-1-5-32-544}). Przed załadowaniem modułów \file{Backup-Create.ps1} i \file{Backup-Restore.ps1} skrypt główny weryfikuje listę ACL i odmawia uruchomienia w przypadku wykrycia praw zapisu dla nieuprzywilejowanych kont.
    \item \textbf{Restrykcje usługi OpenSSH na stacji Windows} - w pliku konfiguracyjnym \file{sshd_config} ograniczono dostęp dyrektywami \code{AllowUsers student@10.0.2.30} oraz \code{PasswordAuthentication no}, a reguła zapory \code{FWBackup-SSH-In} zezwala na ruch przychodzący na port TCP~22 wyłącznie ze wskazanego adresu IP serwera kopii (\code{10.0.2.30}).
    \item \textbf{Separacja uprawnień na serwerze Ubuntu} - skrypt \file{fw-server.sh} należy do konta \code{root} (\code{0755}), katalog kopii \file{/srv/fwbackups/student/backups} posiada uprawnienia \code{2750} (\code{fwbackup:sftpgo}, bit \code{setgid}), a wirtualne konto \code{student} w SFTPGo ma przyznane wyłącznie uprawnienia \code{list} oraz \code{download}.
\end{enumerate}

\section{Weryfikacja funkcjonalna i testy automatyczne}\label{sec:testy-weryfikacja}

Strategię weryfikacji poprawności działania systemu podzielono na trzy uzupełniające się poziomy: zautomatyzowaną analizę statyczną w potoku CI/CD, izolowane testy jednostkowe w piaskownicy (wykonywane bez uprawnień administratora) oraz rzeczywiste testy akceptacyjne na dedykowanej maszynie wirtualnej.

\subsection{Automatyczny potok CI/CD i analiza statyczna kodu}\label{subsec:potok-ci-cd}

W repozytorium kodu wdrożono zautomatyzowany potok CI/CD oparty na GitHub Actions (\file{.github/workflows/ci.yml}), który przy każdym commicie oraz żądaniu pull request uruchamia dwa równoległe zadania weryfikacyjne:
\begin{itemize}
    \item \textbf{Środowisko Windows (\code{windows-latest})} - weryfikuje poprawność składniową kodu PowerShell za pomocą lintera \code{PSScriptAnalyzer} (z rygorystycznym filtrem \code{-Severity Error}), po czym uruchamia testy jednostkowe modułu audytu, interfejsu CLI oraz modułu kopii zapasowych.
    \item \textbf{Środowisko Linux (\code{ubuntu-latest})} - wykonuje analizę statyczną skryptu powłoki serwera \file{fw-server.sh} narzędziem \cmd{shellcheck} (weryfikacja błędów składniowych i podatności w skryptach powłoki), a następnie uruchamia zestaw testów jednostkowych serwera w środowisku \cmd{bash}.
\end{itemize}

Dzięki takiemu podejściu wszelkie regresje składniowe lub naruszenia kontraktów API są natychmiast wychwytywane przed scaleniem zmian.

\subsection{Izolowane testy jednostkowe w piaskownicy (\file{tests/})}\label{subsec:testy-jednostkowe}

Aby umożliwić szybką weryfikację logiki aplikacji na stacjach deweloperskich oraz w kontenerach CI/CD bez konieczności podnoszenia uprawnień do poziomu administratora i bez ingerencji w rzeczywiste reguły sieciowe jądra, opracowano zestaw testów jednostkowych działających w bezpiecznej piaskownicy (\textit{sandbox}):
\begin{enumerate}
    \item \textbf{Audyt transakcyjny i kontrakt AST (\file{Test-FirewallAudit.ps1})} - przekierowuje operacje rejestru do gałęzi \code{HKCU\Software\FwAuditTest-<GUID>}, wykorzystuje lokalny muteks oraz tymczasowy katalog w \file{\%TEMP\%}. Sprawdza integralność łańcucha skrótów SHA-256, procedurę rotacji przy 10~MB oraz politykę \textit{fail-closed}/\textit{fail-open}. Ponadto za pomocą analizatora AST (\code{Language.Parser}) bada cały plik \file{firewall.ps1}, weryfikując, czy wszystkie modyfikacje zapory (cmdlety \cmd{*-NetFirewallRule} i program \cmd{netsh advfirewall}) są hermetycznie opakowane w wywołanie \code{Invoke-AuditedChange}.
    \item \textbf{Interfejs CLI i walidacja danych (\file{Test-CliHelpers.ps1})} - wyodrębnia za pomocą AST funkcje pomocnicze (\code{Read-ProfileChoice}, \code{Read-ProtocolChoice}, \code{Get-AuditPolicyStatusLine}) i testuje je z wykorzystaniem deterministycznej kolejki wejściowej mockującej polecenie \cmd{Read-Host}. Weryfikuje odrzucanie niepoprawnych adresów IPv4/IPv6, masek CIDR i portów spoza zakresu $1\text{-}65535$, automatyczne kierunkowe przypisanie portów oraz format generowanych nazw reguł z prefiksem \code{OSO-}.
    \item \textbf{Klient kopii zapasowych (\file{Test-BackupModule.ps1})} - testuje funkcje modułów \file{Backup-Create.ps1} i \file{Backup-Restore.ps1}. Weryfikuje walidację nazw kont użytkowników wyrażeniem regularnym \code{^[a-z0-9][a-z0-9._-]{0,31}$}, odporność na próby obejścia katalogu (\textit{path traversal}, np.\ \code{../}), obsługę portów SSH/SFTP oraz weryfikację sum kontrolnych SHA-256 pobranych plików polityki \file{.wfw}.
    \item \textbf{Serwer kopii zapasowych Ubuntu (\file{test-fw-server.sh})} - uruchamia wyizolowane procedury skryptu \file{fw-server.sh} w katalogu tymczasowym. Weryfikuje parsowanie pliku konfiguracyjnego \file{clients.conf}, mechanizm blokady współbieżności (\cmd{flock}) oraz politykę retencji archiwum (automatyczne przycinanie najstarszych kopii z zachowaniem maksymalnie 30 ostatnich stanów).
\end{enumerate}

\subsection{Test akceptacyjny end-to-end na maszynie wirtualnej (\file{Invoke-WatcherAcceptance.ps1})}\label{subsec:testy-akceptacyjne}

Weryfikację elementów wymagających rzeczywistych uprawnień administratora oraz interakcji z podsystemami jądra Windows (zadanie harmonogramu \code{SYSTEM}, chroniona lista ACL w \file{C:\ProgramData\FirewallAudit}, odczyt zdarzeń z kanału \code{MPSSVC}) realizuje skrypt akceptacyjny \file{tests\Invoke-WatcherAcceptance.ps1}. Skrypt ten przeprowadza pełny cykl życia systemu na rzeczywistej maszynie wirtualnej:
\begin{enumerate}
    \item weryfikuje ochronę ACL katalogu audytu (brak praw modyfikacji dla grup \code{Users}, \code{Everyone} oraz \code{Authenticated Users}),
    \item instaluje zadanie \code{\FirewallAudit\Watcher} i potwierdza konfigurację 3 wyzwalaczy (zdarzeniowy, startowy, okresowy co 15~min) oraz zgodność skrótu SHA-256 skryptu w katalogu \file{bin},
    \item tworzy i usuwa testową regułę \code{FWAUDIT-ACCEPT-<id>-external} bezpośrednim wywołaniem cmdletu \cmd{New-NetFirewallRule} (poza narzędziem), po czym potwierdza przechwycenie zdarzenia przez zadanie \textit{Watcher} oraz poprawną atrybucję sprawcy z dziennika zdarzeń zapory,
    \item symuluje awarię dziennika poprzez zablokowanie pliku \file{audit.jsonl} wyłącznym uchwytem \code{FileShare.None}, weryfikując zablokowanie zmiany w trybie \code{Closed} oraz prawidłowe odłożenie znacznika i późniejsze uzgodnienie (\textit{reconcile}) w trybie \code{Open},
    \item sprawdza ciągłość łańcucha skrótów SHA-256 i usuwa reguły testowe.
\end{enumerate}

Zbiorcze zestawienie zrealizowanych scenariuszy weryfikacyjnych przedstawiono w tabeli~\ref{tab:wyniki-testow}.

\begin{table}[H]
    \centering
    \caption{Zestawienie scenariuszy weryfikacyjnych i wyników testów wszystkich modułów systemu}\label{tab:wyniki-testow}
    \tablefont
    \begin{tabularx}{\linewidth}{C{1.4cm} L{3.6cm} Y C{1.6cm}}
        \toprule
        \hcell{ID} & \hcell{Badany obszar (Moduł)} & \hcell{Opis scenariusza testowego i kryterium akceptacji} & \hcell{Wynik} \\
        \midrule
        \code{T-01} & Walidacja i numeracja (Moduł 1, \code{WF-01}, \code{WF-02}) & Odrzucanie błędnych adresów IP (ciągi znaków, pojedyncze liczby, oktety $>255$) oraz portów spoza zakresu $1\text{-}65535$; poprawne nadanie numerów \code{05-OSO-*} i \code{06-OSO-*} oraz usuwanie po indeksie. & \cmark \\
        \code{T-02} & Semantyka portów (Moduł 1, \code{WF-01}) & Wymuszenie protokołu TCP/UDP po podaniu portu oraz automatyczne przypisanie \code{LocalPort} dla \textit{Inbound} i \code{RemotePort} dla \textit{Outbound}. & \cmark \\
        \code{T-03} & Profile sieciowe NLA (Moduł 2, \code{WF-03}) & Odczyt aktywnej kategorii interfejsu (\code{Private} na karcie \code{Ethernet}), przypisanie reguł do wybranych profili (\code{Public}, \code{Any}) i ich selektywne egzekwowanie. & \cmark \\
        \code{T-04} & Kontrakt AST i WAL (Moduł 3, \code{WF-04}, \code{WNF-04}) & Analiza statyczna AST pliku \file{firewall.ps1} (brak niezaudytowanych wywołań), trójfazowy zapis \code{intent} $\rightarrow$ \code{change} $\rightarrow$ \code{result}, detekcja \code{scope-escape} i czas migawki $<1$~s. & \cmark \\
        \code{T-05} & Łańcuch skrótów SHA-256 (Moduł 3, \code{WF-06}) & Weryfikacja spójności 77 rekordów w \file{audit.jsonl} oraz wykrycie modyfikacji treści, usunięcia wiersza, obcięcia końca pliku (\textit{head hash}) i podmiany \file{baseline.json}. & \cmark \\
        \code{T-06} & Zadanie \textit{Watcher} (Moduł 3, \code{WF-05}, \code{WNF-02}) & Utworzenie i usunięcie reguły poza narzędziem; przechwycenie zmiany przez zadanie \code{SYSTEM}, atrybucja sprawcy (\code{DESKTOP-M67HB24\Student}) i weryfikacja uprawnień ACL. & \cmark \\
        \code{T-07} & Polityka awarii (Moduł 3, \code{WNF-01}) & Zablokowanie pliku \file{audit.jsonl}: odmowa zmiany w trybie \code{Closed}, wykrycie ręcznej edycji \file{settings.json} oraz uzgodnienie \file{pending.jsonl} w trybie \code{Open}. & \cmark \\
        \code{T-08} & Eksport \textit{pull} (Moduł 4, \code{WF-07}, \code{WNF-03}) & Automatyczne utworzenie kopii na serwerze \code{10.0.2.30} przed dodaniem (\code{auto-add}) i usunięciem reguły (\code{auto-delete}); weryfikacja sumy SHA-256 i grup SFTPGo. & \cmark \\
        \code{T-09} & Odtwarzanie polityki (Moduł 4, \code{WF-08}, \code{WNF-03}) & Przywrócenie kopii \file{.wfw} przez SFTPGo: utworzenie kopii ratunkowej \code{pre-restore}, weryfikacja SHA-256 przed importem i rejestracja \code{Restore} w audycie. & \cmark \\
        \code{T-10} & Analiza statyczna i potok CI/CD (\code{ci.yml}) & Automatyczna weryfikacja reguł jakości kodu w GitHub Actions: \code{PSScriptAnalyzer} (\code{-Severity Error}) dla kodu PowerShell oraz \cmd{shellcheck} dla skryptu \file{fw-server.sh}. & \cmark \\
        \code{T-11} & Testy jednostkowe w piaskownicy (\file{tests/}) & Izolowane testy bez uprawnień administratora: logika wejścia CLI (\file{Test-CliHelpers}), sanityzacja kont i ścieżek (\file{Test-BackupModule}) oraz retencja 30 kopii serwera Linux (\file{test-fw-server.sh}). & \cmark \\
        \bottomrule
    \end{tabularx}
\end{table}
